\section{对单个数据集的描述}
\label{sec:Quality}

CT18 全球分析纳入了 LHC Run-1 的宽范围数据，此外还保留了此前 CT14 分析所用、连同 HERA 合并测量在内的庞大数据库。
Sec.~\ref{sec:data_overview} 与表s.~\ref{tab:EXP_1}--\ref{tab:EXP_2}
综述了 CT18(Z) 数据集，并就每个被拟合实验的 $\chi^2/N_\mathit{pt}$ 与有效高斯
变量 $S_E$ 概括了拟合的总体质量。
然而，要成功拟合全球数据，还必须更细粒度地探究各实验被良好描述的程度。
重要的是用严格的统计度量与检验~\cite{Kovarik:2019xvh} 定量评估数据与理论的符合，包括全面普查拟合各实验时可能出现的张力。我们在 Sec.~\ref{sec:QualityOverview} 勘绘实验约束的全景，主要聚焦拉格朗日乘子（LM）
扫描与敏感度计算这两种互补技术，以阐明拟合内部的符合水平与剩余的系统性张力来源。Section~\ref{sec:Qualitydata} 聚焦于具体被拟合实验的理论描述，而 Section~\ref{sec:EW} 考察 NLO 电弱修正在描述这些数据中的作用。

在流程上，CT 方法的拟合如 App.~\ref{sec:chi2_app} 所述：首先，通过计算与各实验关联系统误差相联系的讨厌参数 $\lambda$ 的最优值，使数据与理论之差极小化；然后，对部分子分布函数函数形式的参数 $a$ 极小化 $\chi^{2}$，得到由式~(\ref{Chi2a0l0}) 给出的最优拟合 $\chi^{2}$，即 $(D^\mathit{sh}_{i}(a_0)-T_{i}(a_0))^{2}/s_{i}^{2}$ 与最优单个讨厌参数 $\overline \lambda (a_0)$ 之平方的和。这里 $T_i$ 是第 $i$ 个理论预言，$D^\mathit{sh}_{i}$ 表示被讨厌参数的最优系统位移平移后的相应数据值；$s_{i}$ 是已发表的总非关联误差估计。

在高质量拟合中，理论偏离数据的程度应与统计和系统不确定度相关的随机涨落一致~\cite{Kovarik:2019xvh}。为核查这一点，我们可以把每个被拟合实验的平移后数据点 $D^\mathit{sh}_{i}$ 与理论值
$T_{i}$ 绘出。
平移后数据的误差棒只含非关联误差
$s_{i}$，因为关联系统误差已由讨厌参数取值计及。

还需考虑第二种比较：绘制与系统不确定度诸源相关联的最优讨厌参数 $\bar \lambda_\alpha(a)$ 的直方图。讨厌参数用于建模可观测量真值与实验测定值之间的关系，通常假定其服从均值 0、标准差 1 的正态分布 ${\cal N}(0,1)$。因此，如果过多的最优拟合参数 $\bar \lambda_\alpha(a)$ 按 ${\cal N}(0,1)$ 远离零，就应当警惕；反之，许多 $\bar \lambda_\alpha(a)$ 接近零——即经验直方图比 ${\cal N}(0,1)$ 更窄——的情形，在若干发表了大量系统不确定度的新数据集中很常见。这种情形一般不那么令人担心，因为出现过多非常小的 $\bar \lambda_\alpha(a)$ 可能有良性原因，参见 \cite{Kovarik:2019xvh} 第 IV.E 节。


\subsection{实验之间的总体一致 \label{sec:QualityOverview}}

\subsubsection{重访有效高斯变量}
先回到图~\ref{fig:sn_ct18}：它基于表s.~\ref{tab:EXP_1} 与 \ref{tab:EXP_2} 收集的信息，展示 CT18 NNLO 全球拟合中各实验单独描述质量的总体情况。
与其考察依赖 $N_{pt,E}$、概率分布各异的 $\chi^2_E(N_{pt,E})/N_{pt,E}$，我们改用等价的信息形式：表s \ref{tab:EXP_1} 与 \ref{tab:EXP_2} 所列有效高斯变量 $S_E=\sqrt{2\chi^2_E}-\sqrt{2N_{pt,E}-1}$ 的直方图~\cite{Lai:2010vv}。

若理论与数据的所有偏离都纯粹来自随机涨落，则无论 $N_{pt,E}$ 多大，$S_E$ 的经验分布都应接近 ${\cal N}(0,1)$。实际上，任何近期的全球拟合给出的 $S_E$
分布在统计上都与 ${\cal N}(0,1)$ 不容~\cite{Kovarik:2019xvh}，表明相对于教科书情形，过多实验欠拟合或过拟合。

对 CT18 NNLO 拟合，图~\ref{fig:sn_ct18} 所示的 $S_E$ 分布最接近 ${\cal N}(0.6,1.9)$。与 ${\cal N}(0,1)$ 相容的概率极小（按 Anderson-Darling 检验~\cite{Kovarik:2019xvh} 为 $p=2.5\cdot 10^{-5}$）。图中我们标记了偏离 $S_E\! =\! 0$ 最大的实验：其一是提供单举 DIS 主要约束的 HERAI+II 合并数据集~\cite{Abramowicz:2015mha}（$S_E\approx 5.7$）——尽管有 Sec.~\ref{sec:summary-HERA2} 讨论的拟合质量问题，全球分析中必须保留它；其二是铁靶带电流 DIS 中结构函数 $x_B F_3(x_B,Q)$ 的 CCFR 测量~\cite{Seligman:1997mc}——中心拟合的 $\chi^2/N_{pt}\approx 0.4$ 异常低，但它确实约束某些味的 PDF 不确定度，例如下一节给出的 LM 扫描所示。

我们还注意到，图~\ref{fig:sn_ct18} 中以浅绿色标出的新 LHC Run-1 数据集，正的 $S_E$ 值多于负值，说明它们的 $\chi^2$ 值大于与已发表实验误差一致的随机涨落预期，查阅表~\ref{tab:EXP_2} 即可证实。

两个方块与两颗星分别标记 NuTeV 双缪子与 CCFR 双缪子数据的 $S_E$ 值；鉴于这些数据对探查奇异性 PDF 的重要性，我们特别予以强调。替代拟合 CT18Z 的类似图（图~\ref{fig:sn_ct18z}）显示，由于 ATLAS 7 TeV $W/Z$ 产生数据的反向拉动，CCFR 与 NuTeV 实验的 $S_E$ 值相对 CT18 拟合有所上升。

\begin{figure}[p]
  \hspace*{-0.6cm}\includegraphics[width=0.45\textwidth]{images/gMHT_scanCT18.pdf}\quad
  \includegraphics[width=0.43\textwidth]{images/gx0p3_scanCT18.pdf}
	\caption{基于 CT18 NNLO 拟合、$Q=125$ GeV 处胶子 PDF 在 $x=0.01$ 与 $0.3$ 的 LM 扫描。
		\label{fig:LMg18}}
\end{figure}


\begin{figure}[p]
	\includegraphics[width=0.49\textwidth]{images/pib23hTnu2E-3p0_0E00_LM27-__1_x2_00E-03_Q1_00E+02_DEchi2_re0_ect.pdf}
	\includegraphics[width=0.49\textwidth]{images/pib23hTnu3E-1p1_0E1_LM27-__1_x3_00E-01_Q1_00E+02_DEchi2_re0_ect.pdf}\\
	\includegraphics[width=0.49\textwidth]{images/pib23hTnd2E-3p0_0E00_LM27-__2_x2_00E-03_Q1_00E+02_DEchi2_re0_ect.pdf}
	\includegraphics[width=0.49\textwidth]{images/pib23hTnd3E-1p0_0E00_LM27-__2_x3_00E-01_Q1_00E+02_DEchi2_re0_ect.pdf}\\ %d_hx.pdf
        	\caption{基于 CT18 拟合、$Q=100$ GeV 处上、下夸克 PDF 在 $x=0.002$ 与 $0.3$ 的 LM 扫描。
		\label{fig:LMud}}
\end{figure}

\begin{figure}[p]
	\center
	\includegraphics[width=0.49\textwidth]{images/pib23hTnub2E-3p0_0E00_LM27-_-1_x2_00E-03_Q1_00E+02_DEchi2_re0_ect.pdf}
	\includegraphics[width=0.49\textwidth]{images/pib23hTnub3E-1p0_0E00_LM27-_-1_x3_00E-01_Q1_00E+02_DEchi2_re0_ect.pdf}\\
	\includegraphics[width=0.49\textwidth]{images/pib23hTndb2E-3p0_0E00_LM27-_-2_x2_00E-03_Q1_00E+02_DEchi2_re0_ect.pdf}
	\includegraphics[width=0.49\textwidth]{images/pib23hTndb3E-1p0_0E00_LM27-_-2_x3_00E-01_Q1_00E+02_DEchi2_re0_ect.pdf}\\
         \includegraphics[width=0.49\textwidth]{images/pib23hTns2E-3p0_0E00_LM27-__3_x2_00E-03_Q1_00E+02_DEchi2_re0_ect.pdf}
	\includegraphics[width=0.49\textwidth]{images/s_hx.pdf}
	\caption{
		同图~\ref{fig:LMud}，此处给出 CT18 中
		$\bar{u}$-、$\bar{d}$- 与 $s$-夸克 PDF 的 LM 扫描。
		}
\label{fig:LMubdbs}
\end{figure}

\subsubsection{拉格朗日乘子扫描
\label{sec:LMScans}
}

文献~\cite{Stump:2001gu} 引入的拉格朗日乘子（LM）扫描技术是评估全球拟合张力水平最稳健的方法之一。该方法借助拉格朗日乘子把某个被拟合分布约束在选定数值上，同时在该约束下重新拟合其余 PDF 参数。这样便可系统地把所选 $x$ 和 $\Q$ 处的 PDF 从无约束全球拟合的偏好值移开。由此产生的 $\chi^2$（或 $S_E$）增量轮廓可对每个被拟合实验逐一计算，揭示 PDF 数值的改动与成功描述具体数据的能力之间的关联程度。

图s.~\ref{fig:LMg18}--\ref{fig:LMduratios} 中的一组面板展示了宽范围 CT18 NNLO PDF 的 LM 扫描 $\chi^2$ 轮廓，标度通常取与高能过程相关的高标度 $\Q\!=\! 100$ GeV，所选取的部分子分数代表低 $x$（$x=0.002$ 与 $0.023$）和高 $x$（$x=0.1$ 与
$0.3$）处的 PDF 行为。扫描的一般特征是：虽然全部实验的全局 $\chi^2$ 在约束良好的 $(x,Q)$ 区域接近抛物线形，某些个别实验可能偏好与全局极小差别相当大的 PDF 值。在全局极小处，这类实验自身的 $\chi^2_E$ 可能被抬高数十个单位。

例如在图~\ref{fig:LMg18} 中，我们给出与胶子密度 $g(x,\Q)$ 相关的两个 LM 扫描。左面板的 LM 探测各最敏感测量对希格斯区胶子 PDF 的拉动——该 PDF 通过占主导的 $gg \to
H$ 道贡献希格斯玻色子产生，尤其在其邻域 $x = m_H /
(14\,\mathrm{TeV})\!\sim\! 0.01$、$\Q\!\sim\!m_H$ 处。显然大部分约束来自 HERA 单举 DIS 数据以及 LHC 喷注数据。

右侧关于 $x=0.3$ 的图中，强约束分散到更多数据集上，特别是高 $p_T$ $Z$ 玻色子 $p_T$ 与顶夸克产生。特别地，ATLAS 7 TeV 单举喷注数据偏好 $g(0.3,125\mbox {GeV})\approx 0.3$，与完整拟合的中心值一致；而 CMS 7 TeV 与 8 TeV 喷注产生分别偏好 $g(0.3,125\mbox {GeV}) = 0.242^{+0.016}_{-0.020}$
与 $0.327^{+0.015}_{-0.010}$ ——按 $\Delta \chi^2=1$ 判据约为 $3\sigma$ 差异。

我们注意到，在某些情形——如图~\ref{fig:LMg18} 右图那样显露出实验间显著张力时——基于动态容忍度的 Hessian 估计~\cite{Martin:2009iq} 可能给出比 LM 扫描总 $\chi^2$ 所示窄得多的 PDF 不确定度，这是相互冲突的实验对 PDF 施加相反拉动之间折中的结果。我们将在 Appendix~\ref{sec:LMRsCT18Z} 进一步讨论，并以图~\ref{fig:rserrors} 给出具体例子。

图~\ref{fig:LMud} 给出 $x=0.002$ 与 $0.3$ 处 $u$、$d$ 夸克 PDF 的 LM 扫描。
在小 $x$ 值处，除 HERA 与 NuTeV 的约束外，LHC $W$ 与 $Z$ 玻色子数据（来自 LHCb、CMS 与 ATLAS）的约束如预期般突出。
在 $x=0.3$，若干固定靶实验（如 CDHSW、BCDMS 与 E866）作出显著贡献。
$d$ 夸克密度以及图~\ref{fig:LMduratios} 所示 $d/u$ 比值的情形类似。

对图~\ref{fig:LMubdbs} 中的 $\bar u$、$\bar d$ 反夸克以及图~\ref{fig:LMduratios} 的 $\bar d/\bar u$ 比值，
从小 $x$ 看，LHCb 数据与 CMS $W$ 玻色子电荷不对称数据扮演重要角色，见图~\ref{fig:LMubdbs}。
另一方面，在大 $x$ 处，味分离依赖 E605、E866 与 NMC 氘核数据。

\begin{figure}[b]
\begin{center}
	\includegraphics[width=0.48\textwidth]{images/Rs_scan2Tct18_1.pdf}\quad
	\includegraphics[width=0.48\textwidth]{images/Rs_scan2Tct18_2.pdf}\\
(a)\hspace{2.6in}(b)\\
	\caption{CT18 NNLO 拟合下 $Q=1.5$ GeV 处对 $R_s$ 的 LM 扫描，$x$ 分别为 0.023 与 0.1。
\label{fig:LMRs}}
\end{center}
\end{figure}

\begin{figure}[tb]
	\center
	\fbox{\parbox[c][3.2cm][c]{0.86\linewidth}{\centering [figure unavailable -- see original paper]\\ \texttt{pib23hTndou2E-3p0\_00E0\_LM27-\_10\_x2\_00E-03\_Q1\_00E+02\_DEchi2\_re0\_ect.pdf}}}
	\fbox{\parbox[c][3.2cm][c]{0.86\linewidth}{\centering [figure unavailable -- see original paper]\\ \texttt{pib23hTndou3E-1p0\_00E0\_LM27-\_10\_x3\_00E-01\_Q1\_00E+02\_DEchi2\_re0\_ect.pdf}}}\\
	\fbox{\parbox[c][3.2cm][c]{0.86\linewidth}{\centering [figure unavailable -- see original paper]\\ \texttt{pib23hTndboub2E-3p0\_00E0\_LM27-\_11\_x2\_00E-03\_Q1\_00E+02\_DEchi2\_re0\_ect.pdf}}}
	\fbox{\parbox[c][3.2cm][c]{0.86\linewidth}{\centering [figure unavailable -- see original paper]\\ \texttt{pjb02a20\_LM25-\_11\_x3\_00E-01\_Q1\_00E+02\_DEchi2\_re0.pdf}}}\\

	\caption{同图~\ref{fig:LMud}，为比值 $d/u$ 与 $\bar d/\bar u$ 的 LM 扫描。
		}
\label{fig:LMduratios}
\end{figure}


LM 方法之威力最直接的展示是图~\ref{fig:LMubdbs} 第三行 CT18 奇异夸克 PDF 的扫描，以及对式~(\ref{eq:Rs}) 定义的奇异性比值 $R_s(x,Q)$ 在 $x=0.023$ 与 $x=0.3$ 的扫描（图~\ref{fig:LMRs}）。从图~\ref{fig:LMubdbs} 左下插图（$x=0.002$）可见，CT18 数据集在 $x < 0.01$ 对 $s(x,Q)$ 无实质直接约束。$s(x,Q)$ 的行为只受到低亮度 ATLAS 7 TeV $W,Z$ 数据（Exp.~ID=268）弱弱的约束，以及探测 $x > 0.01$ 区域的 NuTeV 与 CCFR 双缪子数据向低 $x$ 外推的影响。

在 $x=0.01\!-\!0.1$，图~\ref{fig:LMRs} 中的 $R_s$ 比值显示 NuTeV 与 CCFR 双缪子产生连同 HERA 单举 DIS 的约束占主导，LHCb $W/Z$ 产生及固定靶实验 BCDMS、CDHSW、E866、NMC 的约束较弱。这里扫描揭示一个显著特征：当强迫 $R_s(x,Q)$ 在 $x > 0.01$ 接近 1 时，采用 CT18 奇异性参数化的拟合会变得不稳定。对这些增大了的 $R_s$ 值，$\chi^2$ 值发生波动，或拟合不收敛。$R_s$ 稍大的值在 $x < 0.01$ 尚可容忍。

最后回到图~\ref{fig:LMubdbs} 右下插图中 $x=0.3$ 的 $s(x,Q)$：极大 $x$ 行为再次由较低 $x$（在那里受图中所示诸实验组合约束）的奇异 PDF 外推决定。

从本节可见，拉格朗日乘子途径的优势在于其系统性、稳健性，以及揭示其他技术可能遗漏的张力或不稳定性的能力。
另一方面，该计算需要对 LM 参数的多个取值反复重拟合 PDF——这一局限使 LM 扫描计算代价高昂。

\subsubsection{PDF 敏感度分析
\label{sec:L2}
}
与 Sec.~\ref{sec:LMScans} 探讨的 LM 扫描互补的技术是{\it $L_2$ 敏感度}的计算。$L_2$ 敏感度最早见于文献~\cite{Hobbs:2019gob}，用于分析 CT18(Z) 数据对被拟合 PDF 的拉动的相互作用。这里回顾其基本定义。一种密切相关的实现基于文献~\cite{Wang:2018heo} 详述并在 \texttt{PDFSense} 程序中实现的 $L_1$ 敏感度，将在本节末尾用于给 CT18 数据集的实验按对各种 PDF 组合的敏感度排序。

LM 扫描虽是探索给定分析中被拟合数据集之间可能张力的最稳健手段，但计算代价极高，且只针对特定的 $x$、$Q$ 选择进行。如下所述，$L_2$ 敏感度可以快速计算，并对给定全球分析的 $\Delta \chi^2$ 趋势给出很好的近似。而且 $L_2$ 敏感度可在宽 $x$ 范围内便捷算得，使 LM 扫描所示的 $\Delta \chi^2$ 变化能够同时在多个 $x$ 上可视化并解读。我们强调，下面讨论展示的 $L_2$ 敏感度所得定性结论与基于 LM 扫描本身的图像一致。
尽管对次要实验而言 $L_2$ 敏感度未必总能给出与 LM 扫描相同的数值排序，它们在 LM 扫描未能完全覆盖的更宽 $x$ 范围提供了互补信息。


我们在 Hessian 形式体系~\cite{Pumplin:2002vw,Nadolsky:2008zw,Pumplin:2001ct} 下工作，把每个实验 $E$ 的 $L_2$ 敏感度 $S_{f, L2}(E)$ 计算为

\begin{equation}
S_{f, L2}(E) = \vec{\nabla} \chi^2_E \cdot \frac{ \vec{\nabla} f } { |\vec{\nabla} f| }
             = \Delta \chi^2_E\, \cos \varphi (f, \chi^2_E)\ ,
\label{eq:L2}
\end{equation}
即被拟合 PDF 参数沿 $\vec{\nabla}f$ 方向、离开 $\chi^2(\vec{a})$ 全局极小 $\vec{a}_0$ 作单位长度位移时，对数似然函数 $\chi^2_E$ 的变化。
PDF 参数 $\vec a$ 已归一化，使得任意方向离开最优拟合的单位位移对应 Hessian 误差组的默认置信水平（CT18 为 90\%，平均对应给定方向略小于 $\Delta\chi^2_{\textrm{tot}}=100$）。

该位移使 PDF $f(x,Q)$ 增加其 Hessian PDF 误差 $\Delta f$；只要其 PDF 变化经由关联角
\begin{equation}
	\varphi(f, \chi^2_E) = \cos^{-1} \left( \frac{\vec{\nabla} f}{|\vec{\nabla} f |} \cdot \frac{\vec{\nabla} \chi^2_E}{|\vec{\nabla} \chi^2_E |} \right)\ ,
\end{equation}
与 $\chi^2_E$ 的变化相关联，它就把 $\chi^2_E$ 改变 $\Delta \chi^2_E (\hat{a}_f) = \Delta \chi^2_E\, \cos \varphi (f, \chi^2_E) = S_{f, L2}(E)$。
因此，$L_2$ 敏感度 $S_{f, L2}(E)$ 量化了固定 $x$、$Q$ 处 PDF 的变动对被拟合数据集描述的影响。把 $S_{f, L2}(E)$ 对 $x$ 作图，可获得 CT18(Z) 数据集对全球分析中所拟合 PDF（及其组合）拉动的有用信息。这也便于快速可视化全球拟合内部可能的张力：同一 $(x, Q)$ 处，某味的某些部分子密度的 PDF 变化与 $\chi^2_E$ 的变化正相关（即 $S_{f, L2}(E) > 0$），另一些则反相关（$S_{f, L2}(E) < 0$）。

式~(\ref{eq:L2}) 右端关于 $S_{f,L2}$
的各项计算为
\begin{equation}
\Delta X=\left\vert \vec{\nabla}X\right\vert
=\frac{1}{2}\sqrt{\sum_{i=1}^{N_{\textrm{eig}}}\left(X_{i}^{(+)}-X_{i}^{(-)}\right)^{2}},\label{masterDX}
\end{equation}
与
\begin{equation}
\cos\varphi=\frac{\vec{\nabla}X\cdot\vec{\nabla}Y}{\Delta X\Delta
  Y}=\frac{1}{4\Delta X\,\Delta
  Y}\sum_{i=1}^{N_{\textrm{eig}}}\left(X_{i}^{(+)}-X_{i}^{(-)}\right)\left(Y_{i}^{(+)}-Y_{i}^{(-)}\right),\label{cosphi}
\end{equation}
其中 $X_{i}^{(+)}$ 与 $X_{i}^{(-)}$ 是量 $X$ 沿第 $i$ 个本征矢（$\pm$）方向的参数位移下的取值。用这些对称的主公式，$S_{f,L2}(E)$ 对所有实验求和应在距零几十以内，因为总 $\chi^2$ 的容忍边界近乎球对称。个别实验的 $S_{f,L2}(E)$ 在此精度内彼此趋于抵消；$S_{f,L2}(E)$ 的数量级也可解读为 $E$ 相对其余实验张力的度量。

$L_2$ 敏感度可以对单点残差或最优讨厌参数计算，即式~(\ref{Chi2a0l0}) 的各个部分。一个相关且信息量相近的敏感度定义~\cite{Wang:2018heo} 用残差的绝对值 $|r_i|$ 而非其平方 $r_i^2$ 计算（用 $L_1$ 范数代替 $L_2$ 范数）。

\begin{figure}[!htbp]
	\begin{center}
		\fbox{\parbox[c][3.2cm][c]{0.86\linewidth}{\centering [figure unavailable -- see original paper]\\ \texttt{ifl0\_ct18nn\_L2\_q100\_Sf\_1.pdf}}}
	\end{center}
	\vspace{-2ex}
	\caption{
		对胶子 PDF $g(x,Q\!=\!100\,\mathrm{GeV})$ 拉动最强的 CT18 数据集随 $x$ 变化的 $L_2$ 敏感度。考察某些实验的 $S_{f, L2}(E)$ 在``正方向''达到峰值而另一些实验 sharply 负值的那些 $x$ 区域，可揭示主导数据集间的多处张力。
		例如在 $x\! =\! 0.4$，敏感度曲线显示 CMS 8 TeV 喷注数据（Exp.~ID=545）与 BCDMS $F^d_2$ 数据（Exp.~ID=102）（二者都偏好更大的 $g(0.4,100\mbox{ GeV})$）同 BCDMS $F^p_2$（101）、CDHSW $F_3$（109）、E866 $pp$ Drell-Yan（204）与高 $p_T$ $Z$ 玻色子产生（253）数据集合力向下拉动胶子的强烈竞争。
		在 $x\! \approx \! 0.1$，CMS 8 TeV 喷注数据（Exp.~ID=545）强烈对抗 ATLAS（544）与 CMS 7 TeV（542）喷注产生，以及 CDHSW（108）、CCFR（110）与 BCDMS（102）在各种靶上的 DIS 结构函数 $F_2(x,Q)$ 测量。
	}
\label{fig:L2glu}
\end{figure}

针对 CT18(Z) PDF 及 PDF 比值、反映 CT18 数据集之间相互作用与竞争拉动的大批
$L_2$
敏感度图可在 \cite{CT18L2Sensitivity} 查看。替代拟合 CT18Z 的类似计算见 Sec.~\ref{sec:LMCT18Z}。

图~\ref{fig:L2glu} 展示固定 $Q\! =\! 100$ GeV 处 CT18 数据对胶子 PDF 的 $L_2$ 敏感度，绘出对所有 $x$ 都满足 $|S_{f, L2}(E)| \ge 4$ 的那些实验曲线。
该判据一般识别出所论运动学区域内对 PDF 拉动最强的前 $5-10$ 个实验。
由于邻近希格斯质量标度 $Q\!=\!100$ GeV，图~\ref{fig:L2glu} 突出了若干与 14 TeV 希格斯玻色子产生截面 $\sigma_H (14\,\mathrm{TeV})$ 相关的 CT18 数据集之间的相反拉动，是图~\ref{fig:LMg18}（左）的 $L_2$ 版本。除若干非-LHC 实验（Expt. ID~=101、102、108、109、160、204）在不同 $x$ 处对胶子 PDF 施加显著拉动外，在新拟合的 LHC Run-1 数据中，8 TeV $Z$ $p_T$ ATLAS 数据（Exp.~ID=253）在紧邻 $x=0.01$ 处总体拉动最强：$S_{g, L2}(E)\, \approx\, -(4\!-\!5)$，接近同一邻域内 E866 $pp$ 绝对截面数据（Exp.~ID=204）的拉动。
与此同时，Tevatron 单举（504）与 CMS 8 TeV（545）喷注产生数据的相应拉动在略高的 $x\! = \! 0.05-0.1$ 处更大；在 $x\approx 0.1$，CMS 8 TeV 喷注数据（545）有很强的拉动 $S_{g, L2}(E)\! \approx\! +13$，对抗 ATLAS（544）与 CMS 7 TeV（542）喷注数据集的相反拉动。在更高 $x$，胶子 PDF 的最优拟合行为反映多实验拉动之间的权衡，如图~\ref{fig:L2glu} 题注所详述。
在 $x\! < \! 0.01$，我们注意到 HERA 单举（160）与粲产生（147）数据集之间可见的竞争，其他实验在该区域的约束较不突出。基于图~\ref{fig:L2glu} 的全部观察与我们据 Sec.~\ref{sec:LMScans} LM 扫描的发现一致——以图~\ref{fig:LMg18}（左面板）为代表，那里我们识别出对 $g(x\!=\!0.01, Q\!=\!m_H)$ 施加最严格约束的是同样的实验。

\begin{figure}[p]
\center
  \fbox{\parbox[c][3.2cm][c]{0.86\linewidth}{\centering [figure unavailable -- see original paper]\\ \texttt{ave\_S\_bycolor\_ct18nn.pdf}}}\\
  \fbox{\parbox[c][3.2cm][c]{0.86\linewidth}{\centering [figure unavailable -- see original paper]\\ \texttt{total\_S\_bycolor\_ct18nn.pdf}}}
        	\caption{$L_1$ 敏感度：CT18 NNLO 分析中各实验数据集对各 PDF 味的敏感度，按文献~\cite{Wang:2018heo} 的方法学计算。上（下）插图中格子的颜色（按右方调色板选取）表示纵轴实验集对横轴 PDF 味的点平均（累计）敏感度。
		\label{fig:CT18quilts}}
\end{figure}

本节最后以图~\ref{fig:CT18quilts} 结束：其中给出由 \texttt{PDFSense} 代码按文献~\cite{Wang:2018heo} 的途径计算的 $L_1$
敏感度排序图。
在那篇文章里，我们给出了对 \CTHERAII{} NNLO 分析中实验排序的表格，或按其对 PDF 的总敏感度 $f(x_i,Q_i)$：
\begin{equation}
S_{f,L1}^{\textrm{tot}}(E)\equiv\sum_{i=1}^{N_{pt,E}} \left|S_{f,L1}(i)\right|,
\label{SfL1tot}
\end{equation}
或按每数据点的平均敏感度：
\begin{equation}
S_{f,L1}^{\textrm{ave}}(E)\equiv S_{f,L1}^{\textrm{tot}}(E)/N_{pt,E}.
\label{SfL1ave}
\end{equation}
这两个量分别估计实验 $E$ 在数据点 $i=1,..,N_{pt,E}$ 所探典型 $(x_i, Q_i)$ 处对 PDF $f(x_i,Q_i)$ 的总敏感度（对全部 $N_{pt,E}$ 点求和），或该实验单个数据点的平均敏感度。
两种敏感度便于并排比较各实验约束强度——再次在 Hessian 近似下估计。

图~\ref{fig:CT18quilts} 把文献~\cite{Wang:2018heo} 的排序表图形化，现针对 CT18 NNLO 拟合重新计算，结论大体与 \CTHERAII{} NNLO 所得相似。图~\ref{fig:CT18quilts} 上下两面板分别对应上文讨论的点平均与总敏感度。右侧放置把颜色与 $S_{f,L1}(E)$ 大小联系起来的调色板。
从黄到橙再到红的格子表示（列于左侧的）实验对底部所列 PDF $f(x,\mu)$ 的敏感度递增。
白或灰色格子表示对 $f(x,\mu)$ 敏感度最小的实验。

我们观察到：HERA I+II、BCDMS 与 NMC 数据集的单点敏感度在上面板中相对较低，但当按其大量数据点聚合时，这些实验在下面板中对所有 PDF 味的总敏感度非常大。专用固定靶测量（如 CCFR、NuTeV、E605、E866）正如预期地对特定味（如 $s$、$\bar u$、$\bar d$）最为敏感。

另一方面，若干 LHC 实验具有很强的单点敏感度，尤其是 $t\bar t$ 与高 $p_T$ $Z$ 产生，以及 CMS 7、8 TeV $W$ 电荷不对称（见上面板）。这些实验在下面板中的总敏感度仍相当低，因为其数据点数目少（$N_{pt,E}\approx 10-20$）。相反，ATLAS 与 CMS 7 TeV、尤其 CMS 8 TeV 的单举喷注产生数据集，尽管每数据点敏感度温和，却因数据点多、运动学覆盖广而在所有 LHC 实验中显示出最高的总敏感度。

汇总起来：CT18 拟合中的敏感度主体仍来自 HERA 与固定靶数据，但凭借可观的单点敏感度，LHC 实验已能缩减部分 PDF 不确定度。由于本文前文阐述过的 LHC 实验之间的张力，这些不确定度缩减尚未完全实现。

\subsection{CT18 所拟合数据集的描述}
\label{sec:Qualitydata}
本小节示例 CT18 描述本分析所含单个实验的能力，特别关注新纳入的 LHC Run-1 数据。讨论按具体物理过程组织。


\begin{figure}[tb]
	\fbox{\parbox[c][3.2cm][c]{0.86\linewidth}{\centering [figure unavailable -- see original paper]\\ \texttt{data\_281\_CT18\_abs\_ect.pdf}}}
	\fbox{\parbox[c][3.2cm][c]{0.86\linewidth}{\centering [figure unavailable -- see original paper]\\ \texttt{data\_281\_CT18\_DoT\_v2\_ect.pdf}}}
	\caption{
	D0 Run II 电子电荷不对称数据（Exp.~ID=281）与 CT18 理论的比较。由于不对称在所示范围内过零，右面板显示{\it 差} $\mathrm{Data}\!-\!\mathrm{Theory}$，而非本节他处使用的{\it 比} $\mathrm{Data}/\mathrm{Theory}$。
	}
\label{fig:281}
\end{figure}

\subsubsection{矢量玻色子产生数据}
\label{sec:QualityDYdata}
{\bf Tevatron 电荷不对称。}
CT18 PDF 与固定靶及 Tevatron 实验的矢量玻色子产生数据整体符合良好。特别地，自 CT14~\cite{Dulat:2015mca} 起用于我们分析、对 $x>0.1$ 的 $d(x)/u(x)$ 敏感的高亮度电荷不对称数据集 281（D0 Run-2~\cite{D0:2014kma}）被很好描述。\footnote{根据 $L_2$ 敏感度~\cite{CT18L2Sensitivity}，与全体数据相比，NMC DIS 数据 104 与电荷不对称数据集 281 在大 $x$ 偏好更软的 $d(x)/u(x)$（分别约 15 与 5 个 $\chi^2_E$ 单位）；而 LHCb 7 TeV W 快度（245）与 E866 $pp$ Drell-Yan（204）数据集在同一 $x$ 区偏好更硬的 $d/u$。}

图~\ref{fig:281} 给出电子电荷不对称作为电子赝快度绝对值函数的数据-理论比较。平移后数据用红点表示，未平移数据用黑色。绝对电荷不对称示于图~\ref{fig:281} 左插图；右插图给出 $\mathrm{Data}\!-\!\mathrm{Theory}$ 差，误差棒表示总非关联不确定度（非关联统计与非关联系统误差的平方和），除非另有说明，本文通篇如此一致处理。平移数据相对误差棒的理论偏差体现拟合优度；平移与未平移数据之间的移动则体现关联讨厌参数的效应。
理论预言用 \texttt{ResBos} 代码按近似 NNLO + NNLL 计算。
蓝色带代表用对称 Hessian 误差按 68\% C.L. 评估的 CT18 PDF 不确定度。可以看到数据被 CT18 预言描述良好，仅在一个高赝快度分箱（$|\eta_{e}|\sim2.6$）观察到轻度不符。


\begin{figure}[tbp]
	\fbox{\parbox[c][3.2cm][c]{0.86\linewidth}{\centering [figure unavailable -- see original paper]\\ \texttt{data\_250\_CT18\_\_1\_abs\_ect.pdf}}}
	\fbox{\parbox[c][3.2cm][c]{0.86\linewidth}{\centering [figure unavailable -- see original paper]\\ \texttt{data\_250\_CT18\_\_1\_DoT\_ect.pdf}}}
	\fbox{\parbox[c][3.2cm][c]{0.86\linewidth}{\centering [figure unavailable -- see original paper]\\ \texttt{data\_250\_CT18\_\_2\_abs\_ect.pdf}}}
	\fbox{\parbox[c][3.2cm][c]{0.86\linewidth}{\centering [figure unavailable -- see original paper]\\ \texttt{data\_250\_CT18\_\_2\_DoT\_ect.pdf}}}
	\fbox{\parbox[c][3.2cm][c]{0.86\linewidth}{\centering [figure unavailable -- see original paper]\\ \texttt{data\_250\_CT18\_\_3\_abs\_ect.pdf}}}
	\fbox{\parbox[c][3.2cm][c]{0.86\linewidth}{\centering [figure unavailable -- see original paper]\\ \texttt{data\_250\_CT18\_\_3\_DoT\_ect.pdf}}}
	\caption{LHCb 8 TeV 缪子衰变道 $W^+$（上）、$W^-$（中）、$Z^0$（下）截面测量（Exp.~ID=250）与 CT18 理论预言的比较。数据以每分箱截面 $\sigma=\frac{d\sigma}{d\eta_\mu}\Delta\eta_\mu,\frac{d\sigma}{dy_Z}\Delta y_Z$ 给出，而非微分截面。
	故 $W^-$ 图中直方图分箱 $3.5<\eta_\mu<4.0$ 的凸起源于其较大的分箱宽度。LHCb 7 TeV $W/Z$ 数据（Exp. ID=245）的图中有类似的凸起。
			\label{fig:id250}
}
\end{figure}

\begin{figure}[tb]
	\fbox{\parbox[c][3.2cm][c]{0.86\linewidth}{\centering [figure unavailable -- see original paper]\\ \texttt{data\_245\_CT18\_\_2\_abs\_ect.pdf}}}
	\fbox{\parbox[c][3.2cm][c]{0.86\linewidth}{\centering [figure unavailable -- see original paper]\\ \texttt{data\_245\_CT18\_\_2\_DoT\_ect.pdf}}}
	\fbox{\parbox[c][3.2cm][c]{0.86\linewidth}{\centering [figure unavailable -- see original paper]\\ \texttt{data\_245\_CT18\_\_3\_abs\_ect.pdf}}}
	\fbox{\parbox[c][3.2cm][c]{0.86\linewidth}{\centering [figure unavailable -- see original paper]\\ \texttt{data\_245\_CT18\_\_3\_DoT\_ect.pdf}}}
	\fbox{\parbox[c][3.2cm][c]{0.86\linewidth}{\centering [figure unavailable -- see original paper]\\ \texttt{data\_245\_CT18\_\_1\_abs\_ect.pdf}}}
	\fbox{\parbox[c][3.2cm][c]{0.86\linewidth}{\centering [figure unavailable -- see original paper]\\ \texttt{data\_245\_CT18\_\_1\_DoT\_ect.pdf}}}
	\caption{同图~\ref{fig:id250}，针对 LHCb 7 TeV（Exp.~ID=245）。}
	\label{fig:id245}
\end{figure}


{\bf LHC 数据：LHCb。}
如前所述，在新增的 LHC Drell-Yan 数据集中，LHCb 合作组的 Drell-Yan 截面测量（Exp.~IDs=250、245 与 246，重要性依次递减）对 CT18 PDF 影响最强。与 D0 电子电荷不对称类似，图~\ref{fig:id250} 将 CT18 NNLO 理论预言与 8 TeV 缪子道 $W/Z$ 产生（Exp.~ID=250）的平移和未平移数据都比较。对 7 TeV $\mu$ 道 $W/Z$ 产生（Exp.~ID=245）与 8 TeV $e$ 道 $Z$ 产生（Exp.~ID=246）的类似比较分别见图~\ref{fig:id245} 与图~\ref{fig:LHCb8Zee}。
NNLO 理论用 NLO \texttt{MCFM} 生成的 \texttt{APPLgrid} 文件乘以 \texttt{FEWZ} 与 \texttt{MCFM-8.0} 计算的逐点 $K$ 因子获得。

对于图~\ref{fig:id250} 中 8 TeV 的 $Z/W$ 玻色子产生，理论与数据符合良好，只有快度 2 附近的数据点例外。在所有三个数据集的 $Z$ 玻色子产生中（底行），即便经过系统位移，$y_Z <2.5$ 与 $y_Z=3-4$ 处形状上的理论-数据分歧依然存在。这正是表 \ref{tab:EXP_2} 所引实验 245 与 250 的 $\chi^2_E$ 偏高的原因；纳入 ATLAS 7 TeV $W/Z$ 产生数据集（Expt.~ID=248）后的 CT18Z 拟合部分缓解了该分歧。
$Z$ 产生低快度处，观测轻子运动学接受度存在很大的建模不确定度。
另一方面，$y_Z\approx 4$ 处的分歧与其他全球拟合数据集（例如 CMS 与 D0 $W$ 轻子电荷不对称数据）的拉动存在张力。该张力已用 \texttt{ePump} 程序研究。特别地，我们比较了包含与不包含实验 246 $Z$ 玻色子分布第一（低快度）分箱两种更新拟合。这一选择对最终 PDF 影响很小，不过舍弃第一快度分箱后 LHCb 数据的 $\chi^2_E/N_{pt,E}$ 明显改善。

CT18 拟合对单个数据点的质量可用图~\ref{fig:res_rk_1} 所示平移残差直方图量化。若对实验 $E$ 的拟合良好，其平移残差 $r_i\! =\! (D^\mathit{sh}_{i}-T_{i})/s_{i}$ 与优化讨厌参数 $\bar\lambda_\alpha$ 的直方图应与标准正态分布相符。例如第三幅图展示 LHCb 8 TeV $W^\pm$ 与 $Z$ 数据（Exp.~ID=250）的 $r_i$ 分布，表明 Exp.~ID=250 数据集中有少数数据点取值很大。如预期，大残差来自 $W^\pm$ 与 $Z$ 数据中 $y=2$ 附近的前几个快度分箱，以及 $Z$ 数据中介于 3.5 和 4 之间的 $3.5\lesssim y_Z\lesssim 4$ 区域。
另一个有用的判据是考察拟合 Exp.~ID=250 数据所需的讨厌参数分布，见图~\ref{fig:res_rk_3} 左图。讨厌参数的分布偏离正态分布，有两个讨厌参数取值特别大（$-3$ 与 $+3.7$）。
图~\ref{fig:res_rk_3} 右面板表示这些数据在 $Q=100$ GeV 处对各 PDF 味的 $L_2$ 敏感度。我们看到，相对于完整的全球数据，LHCb 数据偏好 $x<10^{-2}$ 更低的 $u$、$\bar u$ PDF、$x < 10^{-2}$ 略高的 $s$ 以及 $x\approx 0.2$ 更高的 $\bar d$。其他实验与 PDF 组合的 $L_2$ 敏感度图见文献~\cite{CT18L2Sensitivity}。

\begin{figure}[p]
\fbox{\parbox[c][3.2cm][c]{0.86\linewidth}{\centering [figure unavailable -- see original paper]\\ \texttt{data\_246\_CT18\_abs\_ect.pdf}}}
\fbox{\parbox[c][3.2cm][c]{0.86\linewidth}{\centering [figure unavailable -- see original paper]\\ \texttt{data\_246\_CT18\_DoT\_ect.pdf}}}
\caption{LHCb 8 TeV $Z\to e^+e^-$ 产生中 $Z$ 快度分布（Exp.~ID=246）与 CT18 理论预言的比较。
}
\label{fig:LHCb8Zee}
 \end{figure}

\begin{figure}[t]
	\fbox{\parbox[c][3.2cm][c]{0.86\linewidth}{\centering [figure unavailable -- see original paper]\\ \texttt{res\_his\_CT18-245\_4\_ect.pdf}}}
	\fbox{\parbox[c][3.2cm][c]{0.86\linewidth}{\centering [figure unavailable -- see original paper]\\ \texttt{res\_his\_CT18-246\_4\_ect.pdf}}}
	\fbox{\parbox[c][3.2cm][c]{0.86\linewidth}{\centering [figure unavailable -- see original paper]\\ \texttt{res\_his\_CT18-250\_4\_ect.pdf}}}
	\caption{LHCb 7 与 8 TeV $W/Z$ 产生截面的残差分布：Exp.~ID=245（左）、246（中）、250（右）。
		\label{fig:res_rk_1}}
\end{figure}



\begin{figure}[t]
	\fbox{\parbox[c][3.2cm][c]{0.86\linewidth}{\centering [figure unavailable -- see original paper]\\ \texttt{rk\_his\_CT18-250\_\_4\_ect.pdf}}}\quad\fbox{\parbox[c][3.2cm][c]{0.86\linewidth}{\centering [figure unavailable -- see original paper]\\ \texttt{250\_CT18\_L2\_Q100.pdf}}}
	\caption{左：LHCb 8 TeV $W/Z$ 截面（Exp.~ID=250）的讨厌参数分布。\\ 右：这些数据在 $Q=100$ GeV 对 CT18 NNLO PDF 的拉动，按式~(\ref{eq:L2}) 的 $L_2$ 敏感度计算~\cite{CT18L2Sensitivity}。
		\label{fig:res_rk_3}}
\end{figure}



{\bf LHC 数据：CMS 与 ATLAS。}
CMS 合作组 8 TeV 轻子电荷不对称测量（Exp.~ID=249）纳入了所有 CT18 全球拟合。理论预言与平移、未平移数据的比较见图~\ref{fig:id249}。可以看到全部实验数据都在 68\% C.L. PDF 不确定度内拟合良好。



\begin{figure}[t]
	\fbox{\parbox[c][3.2cm][c]{0.86\linewidth}{\centering [figure unavailable -- see original paper]\\ \texttt{data\_249\_CT18\_abs\_ect.pdf}}}
	\fbox{\parbox[c][3.2cm][c]{0.86\linewidth}{\centering [figure unavailable -- see original paper]\\ \texttt{data\_249\_CT18\_DoT\_ect.pdf}}}
	\caption{CMS 8 TeV 电荷不对称数据（Exp.~ID=249）与 CT18 理论预言的比较。}
\label{fig:id249}
\end{figure}


\begin{figure}[p]
\begin{center}
\fbox{\parbox[c][3.2cm][c]{0.86\linewidth}{\centering [figure unavailable -- see original paper]\\ \texttt{CT18-7\_CT18Z-7\_CT14H-7\_1\_mll1\_68CL.pdf}}}
\fbox{\parbox[c][3.2cm][c]{0.86\linewidth}{\centering [figure unavailable -- see original paper]\\ \texttt{CT18-7\_CT18Z-7\_CT14H-7\_2\_mll2\_68CL.pdf}}}
\fbox{\parbox[c][3.2cm][c]{0.86\linewidth}{\centering [figure unavailable -- see original paper]\\ \texttt{CT18-7\_CT18Z-7\_CT14H-7\_3\_mll3\_68CL.pdf}}}
\caption{基于 CT14$_\mathrm{HERAII}$、CT18 与 CT18Z NNLO PDF 的轻子对横向动量分布 $p_{T, \ell \bar \ell}$ 理论预言（QCD 标度取 $\mu_{R}=M_{T, \ell \bar \ell}/2$、$\mu_{F}=M_{T, \ell \bar \ell}$），并与 ATLAS 8 TeV 测量比较。黄色带为对称 Hessian 法计算的 68\% C.L. PDF 不确定度；虚线带为标度不确定度。}
\label{fig:ATL8ZpT}
\end{center}
\end{figure}

\begin{figure}[t]
	\fbox{\parbox[c][3.2cm][c]{0.86\linewidth}{\centering [figure unavailable -- see original paper]\\ \texttt{res\_his\_CT18-253\_4\_ect.pdf}}}
	\fbox{\parbox[c][3.2cm][c]{0.86\linewidth}{\centering [figure unavailable -- see original paper]\\ \texttt{ct18nn\_id253\_lambdas.pdf}}}
	\caption{左：ATLAS 8 TeV $Z$ $p_T$ 数据（Exp.~ID=253）采用式~\ref{scales253_1} 名义 QCD 标度所得的 $\chi^2$ 残差分布。右：$\mu_R=M_{T,\ell\bar\ell}$（默认）与 $M_{T,\ell\bar\ell}/2$ 两种取法（$\mu_F=M_{T,\ell\bar\ell}$）下相应的讨厌参数分布。
		\label{fig:res_rk_6}}
\end{figure}


在 CT18(Z) 分析中，我们还纳入了 ATLAS $\sqrt{s}=8$ TeV 处 $Z$ 衰变轻子对的横向动量（$p_{T}$）分布（Exp.~ID=253）。
这些数据的理论预言基于 $Z+$jet 产生的 NNLO 固定阶计算。
我们强调已施加运动学割断 $45<p_T^Z<150$
GeV，剔除该固定阶计算缺乏必要精度的低、高 $p_T$ 区域。舍弃低 $p_T$ 数据是因为固定阶计算缺失重求和效应。
舍弃高 $p_T$ 数据是因为：（1）考虑到相对较大的统计误差，数据的约束能力有限；（2）电弱修正不可忽略，将在 Sec.~\ref{sec:EW} 讨论。

实际实现时，我们用 \texttt{MCFM} 自制 NLO \texttt{APPLgrid} 文件，再乘以由文献~\cite{Ridder:2015dxa,Gehrmann-DeRidder:2017mvr,Gehrmann-DeRidder:2016jns,Gehrmann-DeRidder:2016zml,Ridder:2016rzm,Ridder:2016nkl} 发表的 NNLO 与 NLO 截面之比构成的 NNLO/NLO $K$ 因子。
为计及 NNLO 理论预言中不可忽略的涨落，我们加入了额外的 0.5\% 理论蒙特卡洛不确定度，由拟合离散 $K$ 因子的光滑曲线的标准差估计。
名义重正化与因子化标度选为

\begin{equation}
\mu_{R}=\mu_{F}=M_{T, \ell \bar \ell}=\sqrt{(p_{T, \ell \bar \ell})^{2}+M_{\ell \bar \ell}^{2}}\; ,
\label{scales253_1}
\end{equation}
假定标度前置因子为 1。
我们还通过把重正化与因子化标度各自独立乘以 2 与 1/2 来研究 QCD 标度依赖性。
具体地，标度不确定度用 7 点标度变化的包络估计：
\begin{equation}
(\mu_{R},\mu_{F}) = \big[(1/2, 1/2),(1, 1/2),(1/2, 1),(1, 1),(1, 2),(2, 1),(2, 2)\big] \times M_{T, \ell \bar \ell}
\;.
\end{equation}

这些 QCD 标度组合都能较好描述 ATLAS $Z p_T$ 数据的形状，但数据偏好高于名义的归一化：既可通过把理论总归一化提高亮度不确定度的 1--2 个标准差来实现，也可减小 $\mu_R$，或将 $\alpha_s$ 提高到 $0.120-0.124$~\cite{Forte:2020pyp}。结果，使用标度
\begin{equation}
\mu_{R}=M_{T, \ell \bar \ell}/2, ~\mu_{F}=M_{T, \ell \bar \ell}\;,
\label{scales253_2}
\end{equation}
ID=253 的 $\chi^2$ 略好，且相对其他标度 PDF 差异可忽略。
采用式~(\ref{scales253_2}) 标度的 CT18 NNLO 理论预言与 ATLAS 8 TeV 数据的比较见图~\ref{fig:ATL8ZpT}。允许总归一化向上移动 $1.2\sigma$ 后，三个不变质量分箱均得到 $\chi^2_E/N_{pt,E}\sim 1$ 且描述良好。
采用式~(\ref{scales253_1}) 标度的补充图（同样 $\chi^2_E/N_{pt,E} \approx 1$，参见表~\ref{tab:EXP_2}，需总归一化移动 $2\sigma$）收录于补充材料。

图~\ref{fig:res_rk_6} 展示这些数据的残差分布（左子图）与讨厌参数分布（右子图）。
可以看到残差分布中理论与数据符合极佳。在讨厌参数分布中，此过程的 101 个讨厌参数里只有一个（与总归一化相关者）在 $\mu_R=M_{T, \ell \bar \ell}/2$（$M_{T, \ell \bar \ell}$）时增加超过 $1\sigma$（$2\sigma$）。七十多个讨厌参数在这些拟合中过于接近零，或许说明实验编目了太多转瞬即逝的系统效应。[这类极小讨厌参数的过剩对 LHC 实验并不罕见，见 Sec.~\ref{sec:Quality} 开头及 \cite{Kovarik:2019xvh} 第 IV.E 节的讨论。]



备择标度 $\mu_{F,R}=M_{\ell \bar \ell}$ 也试过，导致 $p_T$ 分布形状（不只是归一化）描述变差、$\chi^2$ 升高。

剩余差异无法用电弱修正解释，
因为电弱修正又小又负（见 Sec.~\ref{sec:EW}），
会把理论推得更远离数据。相反，归一化的系统位移可能归因于缺失的高阶（N$^{3}$LO）修正，这由两点观察暗示：其一，$Z$ $p_{T}$ 的 NNLO 修正普遍高达 10\%，表明微扰级数收敛缓慢；其二，很大的标度不确定度（约 3-4\%）也提示缺失的高阶效应可能显著。


\subsubsection{喷注数据}
\label{sec:Jet_fit}
历史上，单举喷注产生在约束胶子密度 $g(x,\Q)$ 方面扮演重要角色，Tevatron Run-II 较早喷注数据对 CT10 与 CT14 全球分析的影响即是明证。如今 CT18 还纳入了对撞机能量与亮度更高的 LHC ATLAS 与 CMS 单举喷注产生数据，见 Sec.~\ref{sec:DataJets}。

\begin{figure}[tb]
	\fbox{\parbox[c][3.2cm][c]{0.86\linewidth}{\centering [figure unavailable -- see original paper]\\ \texttt{data\_504\_CT18\_\_com\_DoT\_hori\_ect.pdf}}}
	\caption{CT18 NNLO 与 CDF Run 2 喷注数据（Exp.~ID=504）的 $\mathrm{Data}/\mathrm{Theory}$ 值。
		 \label{fig:id504}}
\end{figure}

\begin{figure}[tb]
	\fbox{\parbox[c][3.2cm][c]{0.86\linewidth}{\centering [figure unavailable -- see original paper]\\ \texttt{data\_514\_CT18\_\_com\_DoT\_hori\_ect.pdf}}}
	\caption{CT18 NNLO 与 D0 Run-2 单举喷注产生（Exp.~ID=514）的 $\mathrm{Data}/\mathrm{Theory}$ 值。
	     \label{fig:id514}}

\end{figure}

{\bf Tevatron Run-II 数据。}
首先考察对 Tevatron Run-II 喷注数据的拟合。
图~\ref{fig:id504} 所示 CDF Run-II
数据未被 NNLO 理论完美描述（按表~\ref{tab:EXP_1} 其 $\chi^2_E/N_{pt,E}\approx 1.7$ 偏高），并且按图~\ref{fig:L2glu} 关于 $g(x,Q)$ 的 $L_2$ 敏感度图，它偏好的胶子 PDF $g(x,Q)$ 形状与所有实验的平均略有不同。图~\ref{fig:id514} 所绘 D0 Run-II 喷注数据与其余数据集符合更好。


{\bf Run-1 LHC 数据。}
在对部分关联系统误差去关联之后（如 Sec.~\ref{sec:DataJets} 与 App.~\ref{sec:ATLASjetdecorrel} 所列），再加上对所有 LHC 喷注数据统一施加的 $0.5\%$ 总体非关联系统误差（见 Sec.~\ref{sec:DataJets} 的讨论），CT18 拟合便能描述图s.~\ref{fig:DoT542}-\ref{fig:DoT545} 所绘的 LHC CMS 与 ATLAS 喷注数据。

虽然 CT18 分析中与理论的符合是合理的，
但我们注意到 LHC 单举喷注数据集彼此之间存在一些张力，尤其是 CMS 7（Exp.~ID=542）与 8 TeV（Exp.~ID=545）的结果。
这些张力在某些运动学区域的特定部分子味上尤为明显——最明显的是胶子 PDF，分别由图s.~\ref{fig:LMg18} 与~\ref{fig:L2glu} 所绘的 LM 扫描和 $L_2$ 敏感度轮廓定量刻画。
对 ATLAS 7 TeV 单举喷注数据，最优拟合要求关联误差把较小快度区的原始数据向下移动，而把高快度原始数据向上移动。补充材料直方图所示 CT18 NNLO PDF 组的大多数最优讨厌参数 $\lambda_{\alpha}$ 分布在 $|\lambda_{\alpha}|\! \sim\! 0$ 附近的狭窄范围。对 CMS 8 TeV 数据集，28 个讨厌参数中有 4 个要求大于 2 的绝对关联移动——比按假定正态统计预期的个数要多。


\begin{figure}[p]
\fbox{\parbox[c][3.2cm][c]{0.86\linewidth}{\centering [figure unavailable -- see original paper]\\ \texttt{data\_542\_CT18\_\_com\_DoT\_hori\_ect.pdf}}}
\caption{CT18 NNLO 与 CMS 7 TeV 单举喷注产生（Exp.~ID=542）的 $\mathrm{Data}/\mathrm{Theory}$ 值。
\label{fig:DoT542}
}
\end{figure}
\begin{figure}[p]
\fbox{\parbox[c][3.2cm][c]{0.86\linewidth}{\centering [figure unavailable -- see original paper]\\ \texttt{data\_544\_CT18\_\_com\_DoT\_hori\_ect.pdf}}}
\caption{CT18 NNLO 与 ATLAS 7 TeV 单举喷注产生（Exp.~ID=544）的 $\mathrm{Data}/\mathrm{Theory}$ 值。
\label{fig:DoT544}
}
\end{figure}
\begin{figure}[tb]
\fbox{\parbox[c][3.2cm][c]{0.86\linewidth}{\centering [figure unavailable -- see original paper]\\ \texttt{data\_545\_CT18\_\_com\_DoT\_hori\_ect.pdf}}}
\caption{CT18 NNLO 与 CMS 8 TeV 单举喷注产生（Exp.~ID=545）的 $\mathrm{Data}/\mathrm{Theory}$ 值。
\label{fig:DoT545}
}
\end{figure}


\subsubsection{顶夸克对产生数据
\label{sec:QualityTopData}
}


CT18 全球分析纳入的两个 $t\bar{t}$ 数据集都被很好地描述，见表~\ref{tab:EXP_2} 后两行给出的 $\chi^2$ 与有效高斯变量 $S_E$ 值。
特别地，对纳入拟合的 CMS（Exp.~ID=573）与 ATLAS（Exp.~ID=580）数据集，
我们分别得到 $S_E\! =\! 0.6$ 与 $S_E\! =\! -1.1$。
为了逐点详尽展示 $t\bar{t}$ 与理论的符合，
图s.~\ref{fig:573} 与~\ref{fig:580} 给出两个数据集的 $(\mathrm{Data})/(\mathrm{Theory})$ 比值图。
为简单起见，图~\ref{fig:573} 中 CMS 数据的误差棒把文献~\cite{Sirunyan:2017azo} 表 5 与表 7 所列统计与关联系统误差一并计入。

这些测量的统计关联（文献~\cite{Sirunyan:2017azo} 表 6）未纳入 CT18 分析，原因是以讨厌参数表示实现该统计关联信息存在技术困难。讨厌参数表示是所有 CT 分析唯一采用的默认表示。CMS 合作组发布的统计关联以协方差矩阵形式给出。
尽管如此，我们用能够处理协方差矩阵形式关联信息的 \texttt{ePumP} 软件独立交叉验证了统计关联的影响。结论是：纳入统计关联对所得 \texttt{ePumP} 更新 PDF 的影响可以忽略。
计入统计关联协方差矩阵后，CMS（Exp.~ID=573）数据的 CT18 基线 $\chi^2$ 值比表 II 所给值（18.9）增大约 5 个单位。（关于把这些测量纳入全球 PDF 分析的相关讨论见文献~\cite{Czakon:2019yrx} 与 \cite{Bailey:2019yze}。）

\begin{figure}[tb]
	\fbox{\parbox[c][3.2cm][c]{0.86\linewidth}{\centering [figure unavailable -- see original paper]\\ \texttt{data\_573\_CT18\_\_com\_DoT\_ect.pdf}}}
	\caption{CMS 8 TeV $t\bar{t}$ 产生数据（Exp.~ID=573）作为顶（反）夸克横动量函数的 $(\mathrm{Data})/(\mathrm{Theory})$ 比较。}
\label{fig:573}
\end{figure}

\begin{figure}[tb]
	\fbox{\parbox[c][3.2cm][c]{0.86\linewidth}{\centering [figure unavailable -- see original paper]\\ \texttt{data\_565\_CT18\_\_1\_DoT.pdf}}}
	\fbox{\parbox[c][3.2cm][c]{0.86\linewidth}{\centering [figure unavailable -- see original paper]\\ \texttt{data\_567\_CT18\_\_2\_DoT.pdf}}}
	\caption{ATLAS 8 TeV $t\bar{t}$ 产生数据（Exp.~ID=580）作为 $t\bar t$ 横向动量（左）与不变质量（右）函数的 $(\mathrm{Data})/(\mathrm{Theory})$ 比较。}
\label{fig:580}
\end{figure}


图~\ref{fig:573} 中，CMS 顶夸克 $p_T$ 分布在本文考察的四个快度分箱内拟合合理。我们发现中间快度分箱 $0.75\! <\! |y_t| \! <\! 0.85$
与 $0.85\! <\! |y_t| \! <\! 0.85$ 之后一段 $0.85\! <\! |y_t| \! <\! 1.45$ 内的一些点上，理论预言与（未）平移数据有适度偏差，造成残差分布稍宽。值得注意的是，如图~\ref{fig:573} 所见，平移（红）与未平移（黑）数据十分相似——拟合 CMS 数据时关联误差的效应相对极小。
尽管如此，关联系统效应对某些截面数值仍然重要，足以让数据移动到与 CT18 预言相差 $1\sigma$ 以内。

相比之下，要非常好地描述图~\ref{fig:580} 所示 ATLAS 类似的
$p_{T,t}$ 与 $m_{t\bar{t}}$ 分布，关键在于使用讨厌参数补偿关联系统效应，见图~\ref{fig:580}。
该数据集大多数分箱的非关联误差很小（小于 1--2 个百分点）。另一方面系统误差可观，系统位移带来理论与数据的极佳符合：CT18 NNLO 的 $\chi^2_E/N_{pt,E}=9.4/15$。


\subsubsection{双缪子产生
\label{sec:Qualitydimuon}
}

中微子深度非弹性散射中的粲夸克产生截面为 $x > 10^{-2}$ 的奇异性 PDF 提供关键的低 $Q$ 约束。
在 CT14 NNLO 分析中，粲夸克产生截面按 QCD NLO~\cite{Gottschalk:1980rv,Gluck:1997sj,Blumlein:2011zu} 在 S-ACOT-$\chi$ 变味数（VFN）方案~\cite{Aivazis:1993kh,Collins:1998rz,Kramer:2000hn,Tung:2001mv} 内计算。
最近，DIS 的带电流系数函数已算至 QCD NNLO，包括夸克质量依赖~\cite{Berger:2016inr,Gao:2017kkx}。
该计算采用 3 味轻夸克的固定味数（FFN）方案，以快速插值表的形式针对 CCFR 与 NuTeV 双缪子实验的运动学发布~\cite{Goncharov:2001qe,Mason:2006qa}。

CT18 分析仍使用 S-ACOT-$\chi$ VFN 方案的 NLO 理论预言，因为它匹配 CCFR 与 NuTeV 实验数据集的精度。
在 NNLO 实现粲夸克质量效应不仅需要 ACOT 型 VFN 方案的 NNLO 带电流截面（目前尚无），还需要与 CCFR、NuTeV 各自系统学研究中所用 QCD 辐射效应实现保持一致。以 NuTeV~\cite{Mason:2006qa} 为例，事件展开、接受度估计、\footnote{CCFR 与 NuTeV 合作组为从双缪子截面提取粲夸克产生截面施加了显著的接受度修正。这些修正仅在 NLO 精度下估计。}以及粲碎片化的研究都用 LO 与 NLO 程序完成，其系统不确定度超过 NNLO 辐射贡献的大小（如文献~\cite{Berger:2016inr,Gao:2017kkx} 所断言），并依赖粲夸克质量和（可以说较小的~\cite{Ball:2018twp}）核修正。CT 分析中，CCFR 与 NuTeV 双缪子截面按 \cite{Mason:2006qa} 第 5.2.1 节的做法假定 $c\to \mu$ 分支比为 0.099。10\% 的归一化不确定度在 $\nu$ 道内视为完全关联，$\bar \nu$ 道同理。其余系统不确定度按平方相加。
总体而言，文献~\cite{Gao:2017kkx} 的讨论表明，对 CCFR 与 NuTeV 的运动学，FFN 方案与任何 VFN 方案 NNLO 结果的差异预计远小于实验数据的精度。

	\begin{figure}[b]
		\begin{center}
			\fbox{\parbox[c][3.2cm][c]{0.86\linewidth}{\centering [figure unavailable -- see original paper]\\ \texttt{CT18dimuonalts1\_KP.pdf}}}\hspace{0.1in}
			\fbox{\parbox[c][3.2cm][c]{0.86\linewidth}{\centering [figure unavailable -- see original paper]\\ \texttt{CT18Zdimuonalts1\_KP.pdf}}}
			\fbox{\parbox[c][3.2cm][c]{0.86\linewidth}{\centering [figure unavailable -- see original paper]\\ \texttt{CT18dimuonalts2\_KP.pdf}}}\hspace{0.1in}
			\fbox{\parbox[c][3.2cm][c]{0.86\linewidth}{\centering [figure unavailable -- see original paper]\\ \texttt{CT18Zdimuonalts2\_KP.pdf}}}
		\end{center}
		\vspace{-2ex}
		\caption{\label{fig:dimuonalt}
			CT18（左）与
			CT18Z（右）NNLO 拟合中的奇异夸克分布 $s(x,Q)$ 与比值 $R_s(x,Q)$，与对 CCFR、NuTeV 测量采用 QCD NNLO 截面的替代拟合比较。
		}
	\end{figure}


作为交叉检验，我们对双缪子产生截面采用 NNLO FFN 计算，进行了记作 CT18(Z)-charmDIS NNLO 的替代拟合。
对 CT18-charmDIS NNLO，全局 $\chi^2$ 减少 6 个单位（相对 CT18），双缪子数据的 $\chi^2_E$ 减少约 1--2 个单位。
对 CT18Z-charmDIS NNLO，相对 CT18Z，全局 $\chi^2$ 与双缪子数据 $\chi^2_E$ 分别减少 11 与 8 个单位。
两种情形中，NNLO 预言与数据的符合都略有改善。

这些选择对奇异夸克 PDF 的影响也作了交叉检验。名义 CT18(Z) 拟合及其``charmDIS NNLO''替代拟合的奇异夸克 PDF $s(x,Q)$ 与式~(\ref{eq:Rs}) 定义的比值 $R_s(x,Q)$ 在图~\ref{fig:dimuonalt} 比较。
在 CT18-charmDIS 拟合中，我们观察到 $x\! \approx\! 0.1$ 处奇异夸克 PDF 略增。这与文献~\cite{Gao:2017kkx} 的 PDF 剖析结果一致，反映同一 $x$ 区负号的 NNLO QCD 修正。在纳入 ATLAS 7 TeV
$W/Z$ 数据的 CT18Z-charmDIS 拟合中，PDF 的变化比 CT18-charmDIS 拟合更小。
此外，双缪子产生 NNLO 贡献所致的变化与 PDF 不确定度的大小相比很小，从名义与替代拟合 $\chi^2$ 值的相对稳定也能推断这一点。


\begin{figure}[htbp]
	\fbox{\parbox[c][3.2cm][c]{0.86\linewidth}{\centering [figure unavailable -- see original paper]\\ \texttt{pdfs\_CT14HERA2dimuons.pdf}}}
	\caption{\label{fig:epump-dimuon}
		不同拟合下 $Q=100$ GeV 处 $s$ PDF 的比较。详见正文。}
\end{figure}

双缪子截面 NNLO 修正使 $x\! \approx\! 0.1$ 处奇异性略微升高的趋势，也已用 \texttt{ePump}~\cite{Hou:2019gfw} 的快速 Hessian 更新技术独立确认，MMHT 组~\cite{Thorne:2019mpt} 亦然，参见 Appendix~\ref{sec:AppendixCT18Z}。把 $c\to \mu$ 分支比从 0.099 改为 MMHT 采用的 0.092~\cite{Harland-Lang:2015nxa} 只会使 CT18 中 $x>0.1$ 的 $s(x,Q)$ 略微升高，同时轻微增大 CCFR+NuTeV 的 $\chi^2$ 值。

最后，为估计粲夸克产生截面的 NNLO 修正对 ATL7ZW 与双缪子数据集同时纳入的影响，
我们完成了一系列 NNLO 拟合，见图~\ref{fig:epump-dimuon}。图中比较四个经 \texttt{ePump} 更新的不同拟合得到的奇异夸克 PDF。PDF 组 (1)
是从 \CTHERAII{} 数据集中去掉 NuTeV 与 CCFR 双缪子数据得到的基线拟合。
把这四个双缪子数据集加回，分别配 NLO 与 NNLO 预言，得到组 (2) 与 (3)。PDF 组 (4)
由加入 ATL7ZW 数据集（不含双缪子数据集）得到。
PDF 组 (3)［用 NNLO 双缪子截面］在 $10^{-3}\! \lesssim\! x\! \lesssim\! 10^{-1}$ 得到的 $s$ PDF 确实略更接近 ATL7ZW 数据约束的结果，但改进仍太弱，不足以化解 ATL7ZW 与双缪子数据集之间的张力。

\subsection{电弱修正}
\label{sec:EW}

\begin{table}
\caption{CT18(Z) 所考虑的 LHC 精密数据的电弱修正小结。对每一过程，我们给出主要观测量、EW 修正的近似上界、计算 EW 修正的参考文献，以及数据是否连同 EW 修正被采纳进 CT18(Z)。}
\label{tab:EWcorrections}
\hspace*{-0.75cm}\begin{tabular}{c|c|c|c|c|c}
\hline
\multirow{2}{*}{Data} & \multirow{2}{*}{Observables}  & Size of EW (and PI)  & Ref. & Data included & EW corrections  \\
 &    & corrections & & in the CT18(Z)? & included in the fits\\
\hline
	Inclusive jet & $p_{T}\sim1.4$ TeV, central  & 8\%  & \cite{Dittmaier:2012kx} & Yes & Yes \\
\hline
	$t\bar{t}$ &$p_{T}^{t}\sim500$ GeV   &  -5\% & \cite{Czakon:2017wor} & Yes & No \\
\hline
$W^{+}(W^{-})$ &     & -0.4(0.3)\%  & \multirow{4}{*}{\cite{Aaboud:2016btc}} & \multirow{4}{*}{CT18Z} & \multirow{4}{*}{Yes}  \\
DY low-mass & $46<M_{\ell \bar \ell}<66$ GeV central  & +1.5\%(PI) +6\%(EW) &  &     \\
DY $Z$-peak  & $66<M_{\ell \bar \ell}<116$ GeV central (forward) & $<$0.1\%(PI)-0.3(-0.4)\%(EW)  & & \\
DY high-mass &$116<M_{\ell \bar \ell}<150$ GeV central (forward) & +1.5\%(PI)-0.5(-1.2)\% (EW) & & \\
\hline
high-mass Drell-Yan & $M_{\ell \bar \ell}\sim1$ TeV  & +5\%(PI)-3\%(EW)  & \texttt{FEWZ} & No  & -- \\
\hline
\multirow{2}{*}{$Z$ $p_{T}$} & $p_{T}\sim m_Z$ & about -5\%   & \cite{Kallweit:2015fta} & Yes & No \\
 & $p_{T}\sim1$ TeV & about -30\%   & \cite{Kallweit:2015fta} & No & --  \\
\hline
\end{tabular}
\end{table}

本小节概述曾为 CT18(Z) 拟合考虑过、并在某些情况下加以应用的电弱（EW）修正。一般而言，我们没有使用 EW 修正很大的数据，尤其当这些数据对 PDF 无显著约束时。EW 修正往往在数据统计误差也很可观的运动学区域更大，因而受 EW 修正影响最大的测量往往对 PDF 较不敏感。注意光子诱导（PI）贡献在受大 EW 修正困扰的运动学区域同样重要，但其符号相反，导致部分抵消。由于 CT18(Z) PDF 未包含显式的光子 PDF，在 EW 修正最大的区域存在高估其影响的风险。对下文所述应用于 CT18(Z) 拟合的那些 EW 修正，实现方式是相乘型 $K$ 因子。

表~\ref{tab:EWcorrections} 总结了 CT18(Z) 所考虑数据 EW 修正的上界，并注明这些数据是否被拟合、EW 修正是否被应用。其中最大者为单举喷注截面的修正，中央快度区最高 $p_{T}$ 分箱可达 8%。
$t\bar{t}$ 产生的 EW 修正已在 Sec.~\ref{sec:TheoryTop} 提及，最大者在 $p_{T}(t)$
分布处。在接近 500 GeV 的高 $p_{T}(t)$ 值，EW 修正为
-5\%，随后在较软 $p_{T}(t)$ 处迅速衰减。
对 $p_{T}(t)$ 谱以外的 $t\bar{t}$ 可观测量，EW 修正相对实验不确定度可忽略。鉴于 8 TeV $t\bar{t}$ 信息在 $p_T\! <\! 500$ GeV 内的实验精度，拟合这些数据时不计入 EW 修正；但描述未来更高 $p_T$ 的测量很可能需要此类修正。

单举 $W^+$、$W^-$ 与 $Z/\gamma^*$
产生数据的 EW 修正已在文献~\cite{Aaboud:2016btc} 用 \texttt{MCSANC} 框架~\cite{Arbuzov:2015yja} 研究。
对 $W^{+}$ 与 $W^{-}$ 产生，发现 EW 修正分别为 $-0.4$\% 与 $-0.3$\%。
在中性流（NC）Drell-Yan（DY）$Z$ 峰区（$66\!<\!M_{\ell \bar \ell}\!<\!116$ GeV）的中央（前向）选择下\footnote{中央选择要求两个轻子都在中央区 $|\eta_l|<2.5$；前向选择要求一个中央轻子和一个前向轻子（$2.5<|\eta_l|<4.9$）。}，EW
修正约 $-0.3(-0.4)\%$，对可观测量 $M_{\ell \bar \ell}$ 与 $y_{\ell \bar \ell}$ 只有微弱的运动学依赖。我们估计光子诱导双轻子产生（$\gamma\gamma\to l^{+}l^{-}$）对 $Z$ 峰 NC DY 的贡献小于 0.1%。对低质量（$46\!<\!M_{\ell \bar \ell}\!<\!66$ GeV）区域，EW 修正为 +6\%，与快度选择判据无关；对高质量（$116\!<\!M_{\ell \bar \ell}\!<\!150$ GeV）NC DY 产生，EW 修正在中央（前向）选择下为 -0.5\%(-1.2\%)，对 $\eta_l$ 与 $y_{\ell \bar \ell}$ 分箱依赖极弱。两个 $M_{\ell \bar \ell}$ 分箱的 PI 贡献均为 1.5\%。
鉴于低、高质量 DY 数据对 PDF 拟合影响甚小，我们决定不在 CT18A(Z) 拟合中纳入低、高质量与前向 $Z$ 峰 DY 数据。
对 $Z$ 峰与 $W^\pm$ 数据，EW 修正包含在相乘 $K$ 因子中，PI 贡献忽略。
最后指出，如文献~\cite{Aaboud:2016btc} 第 6.1.2 节所述，PI 双轻子产生背景已从 ATLAS 7 TeV $W, Z$ 数据中扣除。

出于 EW 修正与 PI 贡献不可忽略的原因，我们也没有把 ATLAS 8 TeV 极高质量（$116\!<\!M_{\ell \bar \ell}\!<\!1500$ GeV）Drell-Yan 数据~\cite{ATLAS8DY} 纳入 CT18(Z) 拟合。我们发现，对极高的不变质量（$M_{\ell \bar \ell}\!\sim\!1$ TeV），用 \texttt{LUXqed17\_plus\_PDF4LHC15}~\cite{Manohar:2017eqh} 计算的 PI 贡献可高达 5%。
相较之下，EW 修正可用 \texttt{FEWZ}
程序计算（见图~\ref{fig:EW4ATL8DY}），此情形约为 -3%。PI 贡献与 EW 修正的部分抵消使截面净增不足 2%。我们还用 \texttt{ePump}
程序核实：这些数据对 CT18 拟合的影响极小。

 \begin{figure}
 \fbox{\parbox[c][3.2cm][c]{0.86\linewidth}{\centering [figure unavailable -- see original paper]\\ \texttt{ATL8DYhiM\_QCDEW\_68CL.pdf}}}
\caption{ATLAS 8 TeV 高质量 Drell-Yan 产生的光子诱导贡献（光子 PDF 取自 \texttt{LUXqed17\_plus\_PDF4LHC15}）与 NLO EW 修正。}
\label{fig:EW4ATL8DY}
 \end{figure}

CT18(Z) 拟合中唯一的 $Z$ $p_{T}$ 分布来自 ATLAS 8 TeV 测量。因缺失的高 $p_T$ 数据 EW 修正显著，我们施加运动学割断 $p_{T}^{Z}\!<\!150$ GeV 舍弃高 $p_T$ 数据。
总体上，这些修正是负的。按文献~\cite{Hollik:2015pja,Kallweit:2015fta}，当
$p_{T}^{Z}\gg M_{Z}$ 时，由于电弱 Sudakov 对数，NLO EW 修正
可达百分之数十。
在被拟合的 $p_{T}^{Z}$ 区间（45 GeV 至 150 GeV），EW 修正使截面减少百分之几，
从而把理论预言进一步推离 ATLAS 8 TeV 数据。
